\subsection{十八道真题的训练定位}

表~\ref{tab:problems} 给出真题与主责能力的对应关系。这里的“主责”表示首先完成
建模和代码基线的人，不表示其他成员不参与复核。每题至少安排一名交叉复核者，
避免模型、数据和论文三条链彼此脱节。

\begin{longtable}{p{1.2cm}p{1.3cm}p{3.3cm}p{6.0cm}p{1.9cm}}
\caption{2020--2025 年 A/B/C 真题训练映射}\label{tab:problems}\\
\toprule
年份 & 题型 & 题目 & 训练任务与模型输入输出关系 & 主责\\
\midrule
\endfirsthead
\toprule
年份 & 题型 & 题目 & 训练任务与模型输入输出关系 & 主责\\
\midrule
\endhead
2020 & A & 炉温曲线 & 建立焊接区中心温度模型并输出给定工况曲线；求最大过炉速度；在制程界限下分别优化高温面积和曲线对称性。 & 队长\\
2020 & B & 穿越沙漠 & 求已知/当日可知天气下的单人策略，并分别研究已知/未知天气下多人资源、路径、挖矿和购买行为的交互策略。 & 组员1\\
2020 & C & 中小微企业信贷策略 & 量化 123 家有记录企业风险并配置固定总额；评估 302 家无记录企业并制定一亿元策略；考虑突发因素调整。 & 组员2\\
2021 & A & FAST 主动反射面形状调节 & 几何关系、受力/位移约束、连续曲面离散、拟合误差和可实现调节量。 & 队长\\
2021 & B & 乙醇偶合制备 C4 烯烃 & 实验因素建模、反应条件优化、多目标权衡及最优条件解释。 & 组员1，组员2复核\\
2021 & C & 生产企业原材料订购与运输 & 供应商评价、订购计划、运输损耗和多周期约束的联立决策。 & 组员1\\
2022 & A & 波浪能装置最大输出功率设计 & 动力学方程、阻尼与刚度参数、数值积分及平均功率寻优。 & 队长\\
2022 & B & 无人机纯方位无源定位 & 方位几何、编队位置误差、迭代校正与多机协同策略。 & 队长，组员1复核\\
2022 & C & 古代玻璃制品成分分析 & 缺失与异常处理、类别差异、成分关系和未知样品分类。 & 组员2\\
2023 & A & 定日镜场优化设计 & 光学/几何效率、遮挡与截断、场地约束及连续—离散联合优化。 & 队长，组员1复核\\
2023 & B & 多波束测线问题 & 海底几何、覆盖宽度、重叠率、测线方向和总航程优化。 & 组员1，队长复核\\
2023 & C & 蔬菜类商品定价与补货决策 & 销量关联、需求预测、损耗、成本加成与日级补货定价。 & 组员2，组员1复核\\
2024 & A & 板凳龙闹元宵 & 螺线与转向几何、碰撞约束、速度传播和队形可行性。 & 队长\\
2024 & B & 生产过程决策问题 & 零配件抽检、装配、拆解、退换损失和多阶段策略比较。 & 组员1\\
2024 & C & 农作物种植策略 & 地块—季节—作物约束、不确定价格产量和多年度轮作计划。 & 组员1，组员2复核\\
2025 & A & 烟幕干扰弹投放策略 & 三维运动、遮蔽几何、时空覆盖和多平台协同投放。 & 队长，组员1复核\\
2025 & B & 碳化硅外延层厚度确定 & 光谱信号处理、物理测量模型、参数反演和误差分析。 & 队长，组员2复核\\
2025 & C & 无创产前检测时点与异常判定 & 分组时点优化、异常识别、风险约束和模型泛化检验。 & 组员2\\
\bottomrule
\end{longtable}

